\import{cardinal/arithmetic.tex}
\import{set/function-space.tex}

\section{Cardinal exponentiation}

\begin{definition}\label{cardinal_two}
    $\cardTwo=2$.
\end{definition}

\begin{proposition}\label{cardinal_two_successor}
    $\cardTwo=\suc{\cardOne}$.
\end{proposition}
\begin{proof}
    Follows by \obv{cardinal_two,cardinal_one}.
\end{proof}

\begin{definition}\label{cardinal_power}
    $\cardPower{\kappa}{\lambda}=\card{\funs{\lambda}{\kappa}}$.
\end{definition}

\begin{proposition}\label{cardinal_two_elements}
    $\cardTwo=\{\cardZero,\cardOne\}$.
\end{proposition}
\begin{proof}
    Show that every element of $\cardTwo$ is an element of
        $\{\cardZero,\cardOne\}$.
    \begin{subproof}
        Fix $x\in\cardTwo$.
        Hence $x=\cardOne$ or $x\in\cardOne$
            by \obv{cardinal_two_successor,suc_elim}.
        Hence $x=\cardOne$ or $x=\cardZero$
            by
            \auto{cardinal_one,cardinal_zero,suc_elim,emptyset}.
        Follows by \auto{cons_iff,singleton_iff}.
    \end{subproof}
    Show that every element of $\{\cardZero,\cardOne\}$ is an element of
        $\cardTwo$.
    \begin{subproof}
        Fix $x\in\{\cardZero,\cardOne\}$.
        Hence $x=\cardZero$ or $x=\cardOne$
            by \auto{cons_iff,singleton_iff}.
        $\cardOne\in\cardTwo$
            by \obv{cardinal_two_successor,suc_intro_self}.
        $\cardZero\in\cardOne$
            by
            \auto{cardinal_zero,cardinal_one,suc_intro_self}.
        Hence $\cardZero\in\cardTwo$
            by \obv{cardinal_two_successor,suc_intro_in}.
    \end{subproof}
    Follows by set extensionality.
\end{proof}

\begin{proposition}\label{cardinal_zero_neq_one}
    $\cardZero\neq\cardOne$.
\end{proposition}
\begin{proof}
    Follows by
        \obv{cardinal_zero,cardinal_one,suc_neq_emptyset}.
\end{proof}

\begin{proposition}\label{powerset_cardinal_one}
    $\pow{\cardOne}=\cardTwo$.
\end{proposition}
\begin{proof}
    Show that every element of $\pow{\cardOne}$ is an element of
        $\cardTwo$.
    \begin{subproof}
        Fix $X\in\pow{\cardOne}$.
        Hence $X\subseteq\cardOne$ by \obv{powerset_elim}.
        \begin{byCase}
            \caseOf{$X=\emptyset$.}
                Hence $X=\cardZero$ by \obv{cardinal_zero}.
                $\cardZero\in\cardTwo$
                    by \auto{cardinal_two_elements,cons_iff,singleton_iff}.
            \caseOf{$X\neq\emptyset$.}
                Hence $X$ is inhabited by \obv{nonempty_inhabited}.
                Take $x$ such that $x\in X$.
                Hence $x\in\cardOne$ by \obv{subseteq}.
                Hence $x=\cardZero$
                    by
                    \auto{cardinal_one,cardinal_zero,suc_elim,emptyset}.
                Hence $\cardZero\in X$.
                Every element of $X$ is equal to $\cardZero$
                    by
                    \auto{subseteq,cardinal_one,cardinal_zero,suc_elim,emptyset}.
                Hence $X=\{\cardZero\}$ by \auto{setext,singleton_iff}.
                Hence $X=\cardOne$
                    by \obv{cardinal_one,cardinal_zero,suc,cons_iff,setext}.
                $\cardOne\in\cardTwo$
                    by \obv{cardinal_two_successor,suc_intro_self}.
        \end{byCase}
    \end{subproof}
    Show that every element of $\cardTwo$ is an element of
        $\pow{\cardOne}$.
    \begin{subproof}
        Fix $X\in\cardTwo$.
        Hence $X=\cardZero$ or $X=\cardOne$
            by \auto{cardinal_two_elements,cons_iff,singleton_iff}.
        $\cardZero\subseteq\cardOne$
            by \obv{cardinal_zero,emptyset_subseteq}.
        $\cardOne\subseteq\cardOne$ by \obv{subseteq_refl}.
        Follows by \auto{powerset_intro}.
    \end{subproof}
    Follows by set extensionality.
\end{proof}

\begin{proposition}\label{cardinal_two_is_cardinal}
    $\cardTwo$ is a cardinal.
\end{proposition}
\begin{proof}
    $\cardOne$ is an ordinal
        by \auto{cardinal_one_is_cardinal,cardinal,initial_ordinal}.
    Hence $\cardTwo$ is an ordinal
        by \obv{cardinal_two_successor,suc_ordinal}.
    Show that for every $\alpha\in\cardTwo$ we have
        $\lnot(\alpha\approx\cardTwo)$.
    \begin{subproof}
        Fix $\alpha\in\cardTwo$.
        Hence $\alpha=\cardZero$ or $\alpha=\cardOne$
            by \auto{cardinal_two_elements,cons_iff,singleton_iff}.
        Suppose $\alpha\approx\cardTwo$.
        \begin{byCase}
            \caseOf{$\alpha=\cardZero$.}
                Hence $\emptyset\approx\cardTwo$
                    by \obv{cardinal_zero}.
                Take $f$ such that
                    $f\in\bijections{\emptyset}{\cardTwo}$
                    by \obv{equinum}.
                $\cardOne\in\cardTwo$
                    by \obv{cardinal_two_successor,suc_intro_self}.
                Take $x\in\emptyset$ such that $f(x)=\cardOne$
                    by \auto{bijections,surj}.
                Contradiction by \obv{emptyset}.
            \caseOf{$\alpha\neq\cardZero$.}
                Hence $\alpha=\cardOne$.
                Hence $\cardOne\approx\pow{\cardOne}$
                    by \obv{powerset_cardinal_one}.
                Take $f$ such that
                    $f\in\bijections{\cardOne}{\pow{\cardOne}}$
                    by \obv{equinum}.
                Hence $f\in\surj{\cardOne}{\pow{\cardOne}}$
                    by \obv{bijections}.
                Contradiction by \auto{cantor}.
        \end{byCase}
    \end{subproof}
    Hence $\cardTwo$ is an initial ordinal by \obv{initial_ordinal}.
    Follows by \obv{cardinal}.
\end{proof}

\begin{proposition}\label{cardinal_power_is_cardinal}
    $\cardPower{\kappa}{\lambda}$ is a cardinal.
\end{proposition}
\begin{proof}
    Follows by \auto{cardinal_power,card_is_cardinal}.
\end{proof}

\begin{proposition}\label{cardinal_power_zero}
    $\cardPower{\kappa}{\cardZero}=\cardOne$.
\end{proposition}
\begin{proof}
    $\funs{\cardZero}{\kappa}=\{\emptyset\}$
        by \auto{cardinal_zero,empty_domain_function_space}.
    $\cardOne=\{\emptyset\}$
        by \obv{cardinal_one,suc,cons_iff,setext}.
    Follows by
        \auto{cardinal_power,cardinal_one_is_cardinal,card_of_cardinal}.
\end{proof}

\begin{proposition}\label{cardinal_power_one_exponent}
    Let $\kappa$ be a cardinal.
    Then $\cardPower{\kappa}{\cardOne}=\kappa$.
\end{proposition}
\begin{proof}
    $\cardOne=\{\emptyset\}$
        by \obv{cardinal_one,suc,cons_iff,setext}.
    Hence $\funs{\cardOne}{\kappa}\approx\kappa$
        by \auto{singleton_domain_function_space_equinum}.
    Hence $\card{\funs{\cardOne}{\kappa}}=\card{\kappa}$
        by \obv{card_eq_iff_equinumerous}.
    $\card{\kappa}=\kappa$ by \obv{card_of_cardinal}.
    Follows by \obv{cardinal_power}.
\end{proof}

\begin{proposition}\label{cardinal_power_one_base}
    $\cardPower{\cardOne}{\lambda}=\cardOne$.
\end{proposition}
\begin{proof}
    $\cardOne=\{\emptyset\}$
        by \obv{cardinal_one,suc,cons_iff,setext}.
    Hence $\funs{\lambda}{\cardOne}
        =\{\lambda\times\{\emptyset\}\}$
        by \auto{singleton_codom_function_space}.
    $\{\lambda\times\{\emptyset\}\}\approx\{\emptyset\}$
        by \obv{singletons_equinum}.
    Hence $\card{\funs{\lambda}{\cardOne}}
        =\card{\{\emptyset\}}$
        by \auto{card_eq_iff_equinumerous}.
    Follows by
        \auto{cardinal_power,cardinal_one_is_cardinal,card_of_cardinal}.
\end{proof}

\begin{theorem}\label{cardinal_power_two_is_powerset_cardinal}
    $\cardPower{\cardTwo}{A}=\card{\pow{A}}$.
\end{theorem}
\begin{proof}
    $\funs{A}{\cardTwo}
        =\funs{A}{\{\cardZero,\cardOne\}}$
        by \obv{cardinal_two_elements}.
    $\pow{A}\approx\funs{A}{\{\cardZero,\cardOne\}}$
        by
        \obv{powerset_equinum_binary_function_space,cardinal_zero_neq_one}.
    Hence $\funs{A}{\cardTwo}\approx\pow{A}$
        by \obv{equinum_sym}.
    Follows by \obv{cardinal_power,card_eq_iff_equinumerous}.
\end{proof}

\begin{theorem}\label{cantor_cardinal_power}
    $\card{A}\in\cardPower{\cardTwo}{A}$.
\end{theorem}
\begin{proof}
    $\cardLt{A}{\pow{A}}$ by \obv{cantor_strict_cardinality}.
    Hence $\card{A}\in\card{\pow{A}}$
        by \obv{card_elem_iff_cardinal_lt}.
    Follows by \obv{cardinal_power_two_is_powerset_cardinal}.
\end{proof}

\begin{corollary}\label{cantor_cardinal_power_for_cardinal}
    Let $\kappa$ be a cardinal.
    Then $\kappa\in\cardPower{\cardTwo}{\kappa}$.
\end{corollary}
\begin{proof}
    $\card{\kappa}=\kappa$ by \obv{card_of_cardinal}.
    Follows by \auto{cantor_cardinal_power}.
\end{proof}

\begin{theorem}\label{cardinal_power_monotone_base}
    Suppose $\cardLe{\kappa}{\mu}$.
    Then $\cardLe{\cardPower{\kappa}{\lambda}}
        {\cardPower{\mu}{\lambda}}$.
\end{theorem}
\begin{proof}
    Take $f$ such that $f\in\inj{\kappa}{\mu}$
        by \auto{cardinal_le}.
    Take $H$ such that
        $H\in\inj{\funs{\lambda}{\kappa}}{\funs{\lambda}{\mu}}$
        by \auto{function_space_codom_injection}.
    Hence $\cardLe{\funs{\lambda}{\kappa}}{\funs{\lambda}{\mu}}$
        by \obv{cardinal_le}.
    Hence $\card{\funs{\lambda}{\kappa}}
        \subseteq\card{\funs{\lambda}{\mu}}$
        by \auto{card_subseteq_iff_cardinal_le}.
    Hence $\cardLe{\card{\funs{\lambda}{\kappa}}}
        {\card{\funs{\lambda}{\mu}}}$
        by \obv{subseteq_cardinal_le}.
    Follows by \obv{cardinal_power}.
\end{proof}

\begin{theorem}\label{cardinal_power_monotone_exponent}
    Suppose $\kappa$ is inhabited and $\cardLe{\lambda}{\mu}$.
    Then $\cardLe{\cardPower{\kappa}{\lambda}}
        {\cardPower{\kappa}{\mu}}$.
\end{theorem}
\begin{proof}
    Take $e$ such that $e\in\inj{\lambda}{\mu}$
        by \auto{cardinal_le}.
    Take $H$ such that
        $H\in\inj{\funs{\lambda}{\kappa}}{\funs{\mu}{\kappa}}$
        by \auto{function_space_domain_injection}.
    Hence $\cardLe{\funs{\lambda}{\kappa}}{\funs{\mu}{\kappa}}$
        by \obv{cardinal_le}.
    Hence $\card{\funs{\lambda}{\kappa}}
        \subseteq\card{\funs{\mu}{\kappa}}$
        by \auto{card_subseteq_iff_cardinal_le}.
    Hence $\cardLe{\card{\funs{\lambda}{\kappa}}}
        {\card{\funs{\mu}{\kappa}}}$
        by \obv{subseteq_cardinal_le}.
    Follows by \obv{cardinal_power}.
\end{proof}

\begin{proposition}\label{cardinal_power_congruent_base}
    Suppose $\kappa\approx\mu$.
    Then $\cardPower{\kappa}{\lambda}
        =\cardPower{\mu}{\lambda}$.
\end{proposition}
\begin{proof}
    $\cardLe{\kappa}{\mu}$ and $\cardLe{\mu}{\kappa}$
        by \obv{equinum_cardinal_le}.
    Hence $\cardLe{\cardPower{\kappa}{\lambda}}
        {\cardPower{\mu}{\lambda}}$ and
        $\cardLe{\cardPower{\mu}{\lambda}}
        {\cardPower{\kappa}{\lambda}}$
        by \obv{cardinal_power_monotone_base}.
    Hence $\cardPower{\kappa}{\lambda}
        \approx\cardPower{\mu}{\lambda}$
        by \obv{cardinal_le_antisymmetric_up_to_equinum}.
    $\cardPower{\kappa}{\lambda}$ and
        $\cardPower{\mu}{\lambda}$ are initial ordinals
        by \auto{cardinal_power_is_cardinal,cardinal}.
    Follows by \obv{equinumerous_initial_ordinals_are_equal}.
\end{proof}

\begin{proposition}\label{cardinal_power_congruent_exponent}
    Suppose $\kappa$ is inhabited and $\lambda\approx\mu$.
    Then $\cardPower{\kappa}{\lambda}
        =\cardPower{\kappa}{\mu}$.
\end{proposition}
\begin{proof}
    $\cardLe{\lambda}{\mu}$ and $\cardLe{\mu}{\lambda}$
        by \obv{equinum_cardinal_le}.
    Hence $\cardLe{\cardPower{\kappa}{\lambda}}
        {\cardPower{\kappa}{\mu}}$ and
        $\cardLe{\cardPower{\kappa}{\mu}}
        {\cardPower{\kappa}{\lambda}}$
        by \obv{cardinal_power_monotone_exponent}.
    Hence $\cardPower{\kappa}{\lambda}
        \approx\cardPower{\kappa}{\mu}$
        by \obv{cardinal_le_antisymmetric_up_to_equinum}.
    $\cardPower{\kappa}{\lambda}$ and
        $\cardPower{\kappa}{\mu}$ are initial ordinals
        by \auto{cardinal_power_is_cardinal,cardinal}.
    Follows by \obv{equinumerous_initial_ordinals_are_equal}.
\end{proof}

\begin{theorem}\label{cardinal_power_product_exponent}
    Suppose $\kappa$ is inhabited.
    Then $\cardPower{\kappa}{\cardProduct{\lambda}{\mu}}
        =\cardPower{\cardPower{\kappa}{\mu}}{\lambda}$.
\end{theorem}
\begin{proof}
    $\cardProduct{\lambda}{\mu}=\card{\lambda\times\mu}$
        by \obv{cardinal_product}.
    $\card{\lambda\times\mu}\approx\lambda\times\mu$
        by \obv{card_equinumerous}.
    Hence $\cardProduct{\lambda}{\mu}\approx\lambda\times\mu$.
    Hence $\cardPower{\kappa}{\cardProduct{\lambda}{\mu}}
        =\cardPower{\kappa}{\lambda\times\mu}$
        by \obv{cardinal_power_congruent_exponent}.
    $\funs{\lambda\times\mu}{\kappa}
        \approx\funs{\lambda}{\funs{\mu}{\kappa}}$
        by \obv{function_space_currying_equinumerous}.
    Hence $\card{\funs{\lambda\times\mu}{\kappa}}
        =\card{\funs{\lambda}{\funs{\mu}{\kappa}}}$
        by \obv{card_eq_iff_equinumerous}.
    Hence $\cardPower{\kappa}{\lambda\times\mu}
        =\cardPower{\funs{\mu}{\kappa}}{\lambda}$
        by \obv{cardinal_power}.
    $\funs{\mu}{\kappa}\approx\card{\funs{\mu}{\kappa}}$
        by \obv{equinumerous_card}.
    Hence $\cardPower{\funs{\mu}{\kappa}}{\lambda}
        =\cardPower{\card{\funs{\mu}{\kappa}}}{\lambda}$
        by \obv{cardinal_power_congruent_base}.
    $\card{\funs{\mu}{\kappa}}=\cardPower{\kappa}{\mu}$
        by \obv{cardinal_power}.
\end{proof}
